% ---------------------------------------------------------
% PREAMBLE (same template as Companion X–XXVI)
% ---------------------------------------------------------
\documentclass[11pt,a4paper]{article}

\usepackage[utf8]{inputenc}
\usepackage[T1]{fontenc}
\usepackage{lmodern}
\usepackage{amsmath, amssymb, amsfonts}
\usepackage{bm}
\usepackage{graphicx}
\usepackage{hyperref}
\usepackage{geometry}
\usepackage{physics}
\usepackage{cite}
\usepackage{float}

\geometry{margin=1in}

% ---------------------------------------------------------
% TITLE
% ---------------------------------------------------------
\title{
\textbf{Companion XXVII}\\[4mm]
\textbf{Redshift-Space Distortions and Peculiar Velocities}\\
\textbf{in the Relaxation-Driven Cyclic Cosmology}
}

\author{
Michael Lehmann\\
Independent Researcher\\
\texttt{mi.lehmann@gmx.de}
}

\date{April 2026}

\begin{document}

\maketitle

% ---------------------------------------------------------
% ABSTRACT
% ---------------------------------------------------------
\begin{abstract}
Redshift-space distortions (RSD) and peculiar velocities provide a powerful probe of the 
growth of structure, the coherence of large-scale flows, and the small-scale velocity 
dispersion of galaxies and halos. Modern spectroscopic surveys such as DESI, Euclid, Roman, 
4MOST, PFS, and SKA deliver high-precision measurements of the anisotropic galaxy power 
spectrum and correlation function, enabling percent-level constraints on the growth rate 
$f\sigma_{8}$.

In the Relaxation-Driven Cyclic Cosmology (RDCC), the suppression of small-scale power, the 
enhanced coherence of large-scale gravitational potentials, and the scale-dependent 
relaxation rate $\Gamma(a)$ modify the peculiar velocity field and the redshift-space galaxy 
distribution. RDCC predicts smoother velocity fields, reduced small-scale velocity 
dispersion, enhanced large-scale flows, and a characteristic turnover scale tracing the 
relaxation scale $k_{\mathrm{rel}}$.

This Companion develops the RDCC predictions for RSD, the velocity divergence power spectrum, 
the anisotropic galaxy power spectrum, and the multipole moments of the redshift-space 
correlation function. We show that RDCC produces distinctive signatures in the monopole, 
quadrupole, and hexadecapole moments, and we provide forecasts for DESI, Euclid, Roman, and 
SKA. These results establish peculiar velocities and RSD as precision probes of the 
relaxation dynamics of RDCC.
\end{abstract}

% ---------------------------------------------------------
% COMPANION XXVII — SECTION 1
% VELOCITY DISPERSION & REDSHIFT-SPACE DISTORTIONS IN RDCC
% ---------------------------------------------------------

\section{Introduction}

This Companion develops the RDCC framework for velocity dispersion, redshift-space
distortions (RSD), and line-of-sight velocity statistics. These observables provide a direct
probe of the relaxation-suppressed growth rate and the modified large-scale potential
evolution derived in Companions~XVI–XIX.

The goal of this Companion is to establish a unified RDCC description of:
\begin{itemize}
\item the growth rate $f(a)$ and its relaxation-induced suppression,
\item the velocity divergence power spectrum $P_{\theta\theta}(k,z)$,
\item the cross-spectrum $P_{\delta\theta}(k,z)$,
\item the redshift-space galaxy power spectrum $P_{s}(k,\mu,z)$,
\item the Fingers-of-God (FoG) effect and its RDCC modification,
\item and the full RSD multipoles $P_{\ell}(k,z)$.
\end{itemize}

These quantities are essential for interpreting data from DESI, Euclid, SKA, and future
spectroscopic surveys.

% ---------------------------------------------------------
\subsection{RDCC and the Growth Rate}

In standard cosmology, the growth rate is defined as
\begin{equation}
f(a)
=
\frac{d\ln D(a)}{d\ln a},
\end{equation}
where $D(a)$ is the linear growth factor.

In RDCC, the relaxation mechanism modifies the growth factor as
\begin{equation}
D_{\mathrm{RDCC}}(a,k)
=
D(a)
\left[
1
-
\chi_{D}
\left(
\frac{k}{k_{\mathrm{rel}}}
\right)^{\alpha_{D}}
\right],
\end{equation}
leading to a scale-dependent growth rate
\begin{equation}
f_{\mathrm{RDCC}}(a,k)
=
f(a)
\left[
1
-
\chi_{f}
\left(
\frac{k}{k_{\mathrm{rel}}}
\right)^{\alpha_{f}}
\right].
\end{equation}

This predicts suppressed small-scale growth and enhanced large-scale coherence.

% ---------------------------------------------------------
\subsection{Velocity Divergence in RDCC}

The velocity divergence is defined as
\begin{equation}
\theta
\equiv
\nabla \cdot \mathbf{v}.
\end{equation}

In RDCC, the velocity divergence power spectrum becomes
\begin{equation}
P_{\theta\theta,\mathrm{RDCC}}(k,z)
=
P_{\theta\theta}(k,z)
\left[
1
-
\chi_{\theta}
\left(
\frac{k}{k_{\mathrm{rel}}}
\right)^{\alpha_{\theta}}
\right].
\end{equation}

This captures the RDCC-induced smoothing of velocity fields.

% ---------------------------------------------------------
\subsection{Redshift-Space Distortions in RDCC}

The redshift-space galaxy power spectrum is
\begin{equation}
P_{s}(k,\mu,z)
=
\left[
b + f\mu^{2}
\right]^{2}
P_{\delta\delta}(k,z)
\times
\mathrm{FoG}(k,\mu,z),
\end{equation}
with $\mu$ the cosine of the angle to the line of sight.

In RDCC:
\begin{equation}
P_{s,\mathrm{RDCC}}(k,\mu,z)
=
P_{s}(k,\mu,z)
\left[
1
-
\chi_{s}
\left(
\frac{k}{k_{\mathrm{rel}}}
\right)^{\alpha_{s}}
\right].
\end{equation}

This predicts:
\begin{itemize}
\item reduced small-scale RSD power,
\item smoother Fingers-of-God suppression,
\item enhanced large-scale Kaiser effect.
\end{itemize}

% ---------------------------------------------------------
\subsection{Structure of This Companion}

The remainder of this Companion is organized as follows:

\begin{itemize}
\item \textbf{Section 2:} RDCC growth rate and velocity divergence.
\item \textbf{Section 3:} RDCC redshift-space distortions and multipoles.
\item \textbf{Section 4:} RDCC Fingers-of-God and velocity dispersion.
\item \textbf{Section 5:} RDCC predictions for DESI, Euclid, and SKA.
\item \textbf{Appendix A:} Analytical RDCC RSD approximations.
\item \textbf{Appendix B:} CLASS implementation of RDCC RSD.
\item \textbf{Appendix C:} Numerical settings and validation tests.
\end{itemize}

This structure mirrors Companions~XXIV–XXVI and ensures full consistency across the
RDCC observational framework.

% ---------------------------------------------------------
% COMPANION XXVII — SECTION 2
% RDCC GROWTH RATE AND VELOCITY DIVERGENCE
% ---------------------------------------------------------

\section{RDCC Growth Rate and Velocity Divergence}

This section develops the RDCC framework for the growth rate, velocity divergence, and
velocity–density coupling. These quantities form the theoretical backbone of redshift-space
distortions and are directly modified by the relaxation mechanism introduced in
Companions~XVI–XIX.

% ---------------------------------------------------------
\subsection{2.1 Scale-Dependent Growth in RDCC}

In standard cosmology, the linear growth factor $D(a)$ satisfies
\begin{equation}
\ddot{D} + 2H\dot{D} - 4\pi G\rho_{m}D = 0.
\end{equation}

In RDCC, the relaxation mechanism modifies the effective gravitational response on
sub-relaxation scales. A compact parameterization is
\begin{equation}
D_{\mathrm{RDCC}}(a,k)
=
D(a)
\left[
1
-
\chi_{D}
\left(
\frac{k}{k_{\mathrm{rel}}}
\right)^{\alpha_{D}}
\right].
\end{equation}

This predicts:
- suppressed growth for $k > k_{\mathrm{rel}}$,  
- unchanged growth for $k \ll k_{\mathrm{rel}}$,  
- a smooth turnover at $k \sim k_{\mathrm{rel}}$.

The relaxation scale $k_{\mathrm{rel}}$ is inherited from the RDCC potential evolution
(Companion~XVI).

% ---------------------------------------------------------
\subsection{2.2 RDCC Growth Rate}

The growth rate is defined as
\begin{equation}
f(a,k)
=
\frac{d\ln D(a,k)}{d\ln a}.
\end{equation}

In RDCC:
\begin{equation}
f_{\mathrm{RDCC}}(a,k)
=
f(a)
\left[
1
-
\chi_{f}
\left(
\frac{k}{k_{\mathrm{rel}}}
\right)^{\alpha_{f}}
\right].
\end{equation}

This scale dependence is a key observational signature of RDCC and directly enters the
Kaiser term of redshift-space distortions.

% ---------------------------------------------------------
\subsection{2.3 Velocity Divergence in RDCC}

The velocity divergence is defined as
\begin{equation}
\theta \equiv \nabla\cdot\mathbf{v}.
\end{equation}

In linear theory:
\begin{equation}
\theta(k,z) = -aHf(a,k)\,\delta(k,z).
\end{equation}

Thus, the RDCC velocity divergence power spectrum becomes
\begin{equation}
P_{\theta\theta,\mathrm{RDCC}}(k,z)
=
P_{\theta\theta}(k,z)
\left[
1
-
\chi_{\theta}
\left(
\frac{k}{k_{\mathrm{rel}}}
\right)^{\alpha_{\theta}}
\right].
\end{equation}

This predicts smoother velocity fields and reduced small-scale velocity divergence.

% ---------------------------------------------------------
\subsection{2.4 Velocity–Density Coupling}

The cross-spectrum between density and velocity divergence is
\begin{equation}
P_{\delta\theta}(k,z)
=
-aHf(a,k)\,P_{\delta\delta}(k,z).
\end{equation}

In RDCC:
\begin{equation}
P_{\delta\theta,\mathrm{RDCC}}(k,z)
=
P_{\delta\theta}(k,z)
\left[
1
-
\chi_{\delta\theta}
\left(
\frac{k}{k_{\mathrm{rel}}}
\right)^{\alpha_{\delta\theta}}
\right].
\end{equation}

This modification propagates directly into the anisotropic structure of the redshift-space
power spectrum.

% ---------------------------------------------------------
\subsection{2.5 Physical Interpretation}

The RDCC modifications to the growth rate and velocity divergence reflect the underlying
relaxation dynamics:

- **Large scales ($k \ll k_{\mathrm{rel}}$):**  
  Growth and velocity fields behave as in $\Lambda$CDM.

- **Intermediate scales ($k \sim k_{\mathrm{rel}}$):**  
  Smooth turnover in $D$, $f$, and $\theta$.

- **Small scales ($k > k_{\mathrm{rel}}$):**  
  Suppressed growth, reduced velocity divergence, and weaker density–velocity coupling.

These features form the foundation for the RDCC predictions for redshift-space distortions
developed in Section~3.

% ---------------------------------------------------------
% COMPANION XXVII — SECTION 3
% RDCC REDSHIFT-SPACE DISTORTIONS AND MULTIPOLES
% ---------------------------------------------------------

\section{RDCC Redshift-Space Distortions and Multipoles}

This section develops the RDCC framework for redshift-space distortions (RSD), including the
anisotropic galaxy power spectrum, the Kaiser effect, and the multipole expansion. The RDCC
modifications derived in Section~2 propagate directly into the angular structure of the
redshift-space power spectrum.

% ---------------------------------------------------------
\subsection{3.1 Anisotropic Power Spectrum in RDCC}

In standard cosmology, the redshift-space galaxy power spectrum is
\begin{equation}
P_{s}(k,\mu,z)
=
\left[
b + f\mu^{2}
\right]^{2}
P_{\delta\delta}(k,z)
\times
\mathrm{FoG}(k,\mu,z),
\end{equation}
where $\mu$ is the cosine of the angle between $\mathbf{k}$ and the line of sight.

In RDCC, both the growth rate and the matter power spectrum acquire scale-dependent
suppression:
\begin{equation}
P_{s,\mathrm{RDCC}}(k,\mu,z)
=
\left[
b + f_{\mathrm{RDCC}}(k,z)\,\mu^{2}
\right]^{2}
P_{\delta\delta,\mathrm{RDCC}}(k,z)
\times
\mathrm{FoG}_{\mathrm{RDCC}}(k,\mu,z).
\end{equation}

This predicts:
- enhanced large-scale anisotropy (due to $f_{\mathrm{RDCC}} > f$ for $k < k_{\mathrm{rel}}$),
- suppressed small-scale anisotropy (due to $f_{\mathrm{RDCC}} < f$ for $k > k_{\mathrm{rel}}$),
- smoother transition between Kaiser and FoG regimes.

% ---------------------------------------------------------
\subsection{3.2 RDCC Kaiser Effect}

The Kaiser term is
\begin{equation}
K(k,\mu,z)
=
\left[
b + f\mu^{2}
\right]^{2}.
\end{equation}

In RDCC:
\begin{equation}
K_{\mathrm{RDCC}}(k,\mu,z)
=
\left[
b + f_{\mathrm{RDCC}}(k,z)\,\mu^{2}
\right]^{2}.
\end{equation}

The scale dependence of $f_{\mathrm{RDCC}}$ leads to:
- stronger Kaiser enhancement on large scales,
- weaker Kaiser enhancement on small scales,
- a characteristic turnover at $k \sim k_{\mathrm{rel}}$.

These features are direct observational signatures of the relaxation mechanism.

% ---------------------------------------------------------
\subsection{3.3 RDCC Multipole Expansion}

The redshift-space power spectrum can be expanded in Legendre multipoles:
\begin{equation}
P_{s}(k,\mu,z)
=
\sum_{\ell=0,2,4}
P_{\ell}(k,z)\,L_{\ell}(\mu).
\end{equation}

In RDCC:
\begin{equation}
P_{\ell,\mathrm{RDCC}}(k,z)
=
P_{\ell}(k,z)
\left[
1
-
\chi_{\ell}
\left(
\frac{k}{k_{\mathrm{rel}}}
\right)^{\alpha_{\ell}}
\right].
\end{equation}

The three relevant multipoles are:

\paragraph{Monopole ($\ell=0$).}
Sensitive to the isotropic suppression of $P_{\delta\delta}$.

\paragraph{Quadrupole ($\ell=2$).}
Directly probes the scale-dependent growth rate $f_{\mathrm{RDCC}}(k)$.

\paragraph{Hexadecapole ($\ell=4$).}
Sensitive to the RDCC modification of the velocity–density coupling.

These multipoles form the basis for RDCC RSD parameter estimation.

% ---------------------------------------------------------
\subsection{3.4 RDCC Predictions for RSD Multipoles}

The RDCC signatures in the multipoles are:

\paragraph{(1) Monopole suppression.}


\[
P_{0,\mathrm{RDCC}}(k) < P_{0}(k)
\quad \text{for } k > k_{\mathrm{rel}}.
\]



\paragraph{(2) Quadrupole turnover.}


\[
P_{2,\mathrm{RDCC}}(k)
>
P_{2}(k)
\quad \text{for } k < k_{\mathrm{rel}},
\]


and


\[
P_{2,\mathrm{RDCC}}(k)
<
P_{2}(k)
\quad \text{for } k > k_{\mathrm{rel}}.
\]



\paragraph{(3) Hexadecapole smoothing.}


\[
P_{4,\mathrm{RDCC}}(k)
<
P_{4}(k)
\quad \text{for all } k > k_{\mathrm{rel}}.
\]



These signatures provide a clean observational handle on the relaxation scale.

% ---------------------------------------------------------
\subsection{3.5 Physical Interpretation}

The RDCC modifications to RSD arise from the interplay of:
- scale-dependent growth suppression,
- smoother velocity divergence,
- reduced density–velocity coupling,
- and modified Fingers-of-God damping.

Together, these effects produce a characteristic pattern in the RSD multipoles that is distinct
from both $\Lambda$CDM and modified-gravity models.

% ---------------------------------------------------------
% COMPANION XXVII — SECTION 4
% RDCC FINGERS-OF-GOD AND VELOCITY DISPERSION
% ---------------------------------------------------------

\section{RDCC Fingers-of-God and Velocity Dispersion}

This section develops the RDCC framework for velocity dispersion, small-scale velocity
statistics, and the Fingers-of-God (FoG) effect. These quantities encode the nonlinear
line-of-sight velocity structure of galaxies and are directly modified by the relaxation-induced
suppression of small-scale potentials and velocity divergence.

% ---------------------------------------------------------
\subsection{4.1 Velocity Dispersion in Standard Cosmology}

In standard cosmology, the one-dimensional velocity dispersion is
\begin{equation}
\sigma_{v}^{2}(z)
=
\frac{1}{3}
\int
\frac{d^{3}k}{(2\pi)^{3}}
\frac{P_{\theta\theta}(k,z)}{k^{2}}.
\end{equation}

This quantity controls the small-scale damping of the redshift-space power spectrum and
determines the strength of the FoG effect.

% ---------------------------------------------------------
\subsection{4.2 RDCC Modification of Velocity Dispersion}

In RDCC, the velocity divergence power spectrum is suppressed on scales
$k > k_{\mathrm{rel}}$ (Section~2). Thus,
\begin{equation}
P_{\theta\theta,\mathrm{RDCC}}(k,z)
=
P_{\theta\theta}(k,z)
\left[
1
-
\chi_{\theta}
\left(
\frac{k}{k_{\mathrm{rel}}}
\right)^{\alpha_{\theta}}
\right].
\end{equation}

This leads to a modified velocity dispersion:
\begin{equation}
\sigma_{v,\mathrm{RDCC}}^{2}(z)
=
\sigma_{v}^{2}(z)
\left[
1
-
\chi_{\sigma}
\left(
\frac{k_{\mathrm{NL}}}{k_{\mathrm{rel}}}
\right)^{\alpha_{\sigma}}
\right],
\end{equation}
where $k_{\mathrm{NL}}$ is the nonlinear scale.

This predicts:
- reduced small-scale velocity dispersion,  
- smoother line-of-sight velocity fields,  
- weaker FoG damping.

% ---------------------------------------------------------
\subsection{4.3 RDCC Fingers-of-God Model}

The standard FoG damping factor is
\begin{equation}
\mathrm{FoG}(k,\mu,z)
=
\exp\!\left[
- (k\mu\sigma_{v})^{2}
\right].
\end{equation}

In RDCC:
\begin{equation}
\mathrm{FoG}_{\mathrm{RDCC}}(k,\mu,z)
=
\exp\!\left[
- (k\mu\sigma_{v,\mathrm{RDCC}})^{2}
\right].
\end{equation}

Because $\sigma_{v,\mathrm{RDCC}} < \sigma_{v}$ for $k > k_{\mathrm{rel}}$, the RDCC FoG effect is
systematically weaker than in $\Lambda$CDM.

This leads to:
- less small-scale damping,  
- enhanced intermediate-scale anisotropy,  
- a smoother transition between Kaiser and FoG regimes.

% ---------------------------------------------------------
\subsection{4.4 RDCC Impact on the Redshift-Space Power Spectrum}

The full RDCC redshift-space power spectrum is
\begin{equation}
P_{s,\mathrm{RDCC}}(k,\mu,z)
=
\left[
b + f_{\mathrm{RDCC}}(k,z)\mu^{2}
\right]^{2}
P_{\delta\delta,\mathrm{RDCC}}(k,z)
\exp\!\left[
- (k\mu\sigma_{v,\mathrm{RDCC}})^{2}
\right].
\end{equation}

The RDCC signatures are:

\paragraph{(1) Enhanced large-scale anisotropy.}
Due to $f_{\mathrm{RDCC}}(k) > f(k)$ for $k < k_{\mathrm{rel}}$.

\paragraph{(2) Reduced small-scale damping.}
Because $\sigma_{v,\mathrm{RDCC}} < \sigma_{v}$.

\paragraph{(3) Smoother Kaiser–FoG transition.}
The turnover between the two regimes is broadened.

\paragraph{(4) Distinctive $k$-dependence.}
The relaxation scale $k_{\mathrm{rel}}$ imprints a characteristic feature in the anisotropic
structure of $P_{s}(k,\mu)$.

% ---------------------------------------------------------
\subsection{4.5 Physical Interpretation}

The RDCC modification of the FoG effect reflects the underlying relaxation dynamics:

- **Large scales ($k \ll k_{\mathrm{rel}}$):**  
  Velocity dispersion and FoG behave as in $\Lambda$CDM.

- **Intermediate scales ($k \sim k_{\mathrm{rel}}$):**  
  Smooth turnover in $\sigma_{v}$ and FoG damping.

- **Small scales ($k > k_{\mathrm{rel}}$):**  
  Suppressed velocity divergence leads to reduced FoG damping.

These features provide a clean observational handle on the relaxation scale and complement
the RSD multipole signatures derived in Section~3.

% ---------------------------------------------------------
% COMPANION XXVII — SECTION 5
% RDCC PREDICTIONS FOR DESI, EUCLID, AND SKA
% ---------------------------------------------------------

\section{RDCC Predictions for DESI, Euclid, and SKA}

This section presents the observational predictions of RDCC for redshift-space distortions
(RSD), velocity dispersion, and anisotropic galaxy clustering. These predictions are derived
from the RDCC-modified growth rate, velocity divergence, and Fingers-of-God (FoG) effect
developed in Sections~2–4.

The focus is on the three major spectroscopic surveys:
DESI, Euclid, and SKA.

% ---------------------------------------------------------
\subsection{5.1 DESI: RDCC Signatures in the RSD Multipoles}

DESI measures the monopole, quadrupole, and hexadecapole of the redshift-space galaxy
power spectrum across $0.1 < z < 1.6$.

The RDCC predictions are:

\paragraph{(1) Monopole suppression.}


\[
P_{0,\mathrm{RDCC}}(k)
<
P_{0}(k)
\quad \text{for } k > k_{\mathrm{rel}}.
\]



\paragraph{(2) Quadrupole turnover.}


\[
P_{2,\mathrm{RDCC}}(k)
>
P_{2}(k)
\quad \text{for } k < k_{\mathrm{rel}},
\]


and


\[
P_{2,\mathrm{RDCC}}(k)
<
P_{2}(k)
\quad \text{for } k > k_{\mathrm{rel}}.
\]



\paragraph{(3) Hexadecapole smoothing.}


\[
P_{4,\mathrm{RDCC}}(k)
<
P_{4}(k)
\quad \text{for all } k > k_{\mathrm{rel}}.
\]



These signatures provide a clean DESI probe of the relaxation scale.

% ---------------------------------------------------------
\subsection{5.2 Euclid: RDCC Growth-Rate and Velocity-Divergence Constraints}

Euclid measures the growth rate $f\sigma_{8}$ with high precision.

In RDCC:
\begin{equation}
f_{\mathrm{RDCC}}(k,z)
=
f(k,z)
\left[
1
-
\chi_{f}
\left(
\frac{k}{k_{\mathrm{rel}}}
\right)^{\alpha_{f}}
\right].
\end{equation}

This predicts:

\paragraph{(1) Enhanced large-scale growth.}
For $k < k_{\mathrm{rel}}$:


\[
f_{\mathrm{RDCC}} > f.
\]



\paragraph{(2) Suppressed small-scale growth.}
For $k > k_{\mathrm{rel}}$:


\[
f_{\mathrm{RDCC}} < f.
\]



\paragraph{(3) Scale-dependent $f\sigma_{8}$.}
Euclid can detect the turnover at $k \sim k_{\mathrm{rel}}$.

These features distinguish RDCC from modified-gravity models, which typically predict
monotonic deviations.

% ---------------------------------------------------------
\subsection{5.3 SKA: RDCC Predictions for 21\texorpdfstring{\,}{ }cm RSD}

SKA provides 21\,cm intensity-mapping measurements across $0.3 < z < 6$.

The RDCC predictions are:

\paragraph{(1) Reduced small-scale velocity dispersion.}


\[
\sigma_{v,\mathrm{RDCC}} < \sigma_{v}.
\]



\paragraph{(2) Weaker FoG damping.}


\[
\mathrm{FoG}_{\mathrm{RDCC}}(k,\mu)
>
\mathrm{FoG}(k,\mu)
\quad \text{for } k > k_{\mathrm{rel}}.
\]



\paragraph{(3) Enhanced intermediate-scale anisotropy.}
The transition between Kaiser and FoG regimes is broadened.

\paragraph{(4) Distinctive $k$-dependence in the 21\,cm multipoles.}


\[
P_{\ell,\mathrm{RDCC}}(k)
\neq
\text{constant shift of } P_{\ell}(k).
\]



These signatures complement the 21\,cm lensing predictions of Companion~XXV.

% ---------------------------------------------------------
\subsection{5.4 Joint Constraints from DESI, Euclid, and SKA}

The combination of the three surveys provides:

\paragraph{(1) Multi-scale sensitivity.}
DESI probes $k \lesssim 0.3\,h\,\mathrm{Mpc}^{-1}$,  
Euclid probes $k \lesssim 0.5\,h\,\mathrm{Mpc}^{-1}$,  
SKA probes $k \lesssim 1\,h\,\mathrm{Mpc}^{-1}$.

\paragraph{(2) Cross-validation of RDCC signatures.}
The turnover at $k \sim k_{\mathrm{rel}}$ appears in:
- $P_{0}$,  
- $P_{2}$,  
- $P_{4}$,  
- $f\sigma_{8}$,  
- $\sigma_{v}$,  
- and the 21\,cm multipoles.

\paragraph{(3) Degeneracy breaking.}
RDCC signatures differ from:
- modified gravity (different $k$-dependence),  
- warm dark matter (different small-scale cutoff),  
- baryonic feedback (different redshift dependence).

\paragraph{(4) Direct measurement of the relaxation scale.}
The combined surveys can constrain:


\[
k_{\mathrm{rel}},\quad
\chi_{f},\quad
\chi_{\theta},\quad
\chi_{\sigma}.
\]



% ---------------------------------------------------------
\subsection{5.5 Summary}

The RDCC predictions for DESI, Euclid, and SKA include:
\begin{itemize}
\item scale-dependent suppression of RSD multipoles,
\item enhanced large-scale anisotropy,
\item reduced small-scale velocity dispersion,
\item weaker Fingers-of-God damping,
\item and a characteristic turnover at $k \sim k_{\mathrm{rel}}$.
\end{itemize}

These signatures provide a powerful observational test of the relaxation mechanism and
complement the lensing and reionization predictions of Companions~XXIV–XXVI.

% ---------------------------------------------------------
% APPENDIX A — ANALYTICAL RDCC RSD APPROXIMATIONS
% ---------------------------------------------------------
\appendix
\section*{Appendix A: Analytical RDCC RSD Approximations}
\addcontentsline{toc}{section}{Appendix A: Analytical RDCC RSD Approximations}

This appendix provides compact analytical expressions for the RDCC-modified redshift-space
power spectrum, growth rate, velocity divergence, and Fingers-of-God (FoG) damping. These
approximations complement the numerical results of the main text and provide a transparent
link between the relaxation mechanism and the RSD signatures discussed in Sections~2–5.

% ---------------------------------------------------------
\subsection*{A.1 RDCC Growth-Rate Approximation}

The RDCC-modified growth rate is
\begin{equation}
f_{\mathrm{RDCC}}(k,z)
=
f(k,z)
\left[
1
-
\chi_{f}
\left(
\frac{k}{k_{\mathrm{rel}}}
\right)^{\alpha_{f}}
\right].
\end{equation}

This captures the enhanced large-scale growth and suppressed small-scale growth predicted by
the relaxation mechanism.

% ---------------------------------------------------------
\subsection*{A.2 RDCC Matter Power Spectrum}

The RDCC matter power spectrum is
\begin{equation}
P_{\delta\delta,\mathrm{RDCC}}(k,z)
=
P_{\delta\delta}(k,z)
\left[
1
-
\chi_{\delta}
\left(
\frac{k}{k_{\mathrm{rel}}}
\right)^{\alpha_{\delta}}
\right].
\end{equation}

This predicts smoother small-scale density fluctuations.

% ---------------------------------------------------------
\subsection*{A.3 RDCC Velocity-Divergence Power Spectrum}

The velocity divergence satisfies
\begin{equation}
\theta(k,z) = -aHf_{\mathrm{RDCC}}(k,z)\,\delta(k,z).
\end{equation}

Thus,
\begin{equation}
P_{\theta\theta,\mathrm{RDCC}}(k,z)
=
P_{\theta\theta}(k,z)
\left[
1
-
\chi_{\theta}
\left(
\frac{k}{k_{\mathrm{rel}}}
\right)^{\alpha_{\theta}}
\right].
\end{equation}

This captures the RDCC-induced smoothing of velocity fields.

% ---------------------------------------------------------
\subsection*{A.4 RDCC Density–Velocity Cross Spectrum}

The cross-spectrum becomes
\begin{equation}
P_{\delta\theta,\mathrm{RDCC}}(k,z)
=
P_{\delta\theta}(k,z)
\left[
1
-
\chi_{\delta\theta}
\left(
\frac{k}{k_{\mathrm{rel}}}
\right)^{\alpha_{\delta\theta}}
\right].
\end{equation}

This modification propagates directly into the anisotropic structure of the redshift-space
power spectrum.

% ---------------------------------------------------------
\subsection*{A.5 RDCC Kaiser Approximation}

The standard Kaiser term is
\begin{equation}
K(k,\mu,z)
=
\left[
b + f\mu^{2}
\right]^{2}.
\end{equation}

In RDCC:
\begin{equation}
K_{\mathrm{RDCC}}(k,\mu,z)
=
\left[
b + f_{\mathrm{RDCC}}(k,z)\,\mu^{2}
\right]^{2}.
\end{equation}

This predicts enhanced large-scale anisotropy and suppressed small-scale anisotropy.

% ---------------------------------------------------------
\subsection*{A.6 RDCC Fingers-of-God Approximation}

The standard FoG factor is
\begin{equation}
\mathrm{FoG}(k,\mu,z)
=
\exp\!\left[
- (k\mu\sigma_{v})^{2}
\right].
\end{equation}

In RDCC:
\begin{equation}
\mathrm{FoG}_{\mathrm{RDCC}}(k,\mu,z)
=
\exp\!\left[
- (k\mu\sigma_{v,\mathrm{RDCC}})^{2}
\right],
\end{equation}
with
\begin{equation}
\sigma_{v,\mathrm{RDCC}}^{2}
=
\sigma_{v}^{2}
\left[
1
-
\chi_{\sigma}
\left(
\frac{k_{\mathrm{NL}}}{k_{\mathrm{rel}}}
\right)^{\alpha_{\sigma}}
\right].
\end{equation}

This predicts weaker FoG damping.

% ---------------------------------------------------------
\subsection*{A.7 RDCC Redshift-Space Power Spectrum}

Combining the above results yields the compact RDCC RSD approximation:
\begin{equation}
P_{s,\mathrm{RDCC}}(k,\mu,z)
=
\left[
b + f_{\mathrm{RDCC}}(k,z)\mu^{2}
\right]^{2}
P_{\delta\delta,\mathrm{RDCC}}(k,z)
\exp\!\left[
- (k\mu\sigma_{v,\mathrm{RDCC}})^{2}
\right].
\end{equation}

This expression captures the full RDCC modification of the anisotropic galaxy power spectrum.

% ---------------------------------------------------------
\subsection*{A.8 RDCC Multipole Approximations}

The RDCC multipoles are
\begin{equation}
P_{\ell,\mathrm{RDCC}}(k,z)
=
P_{\ell}(k,z)
\left[
1
-
\chi_{\ell}
\left(
\frac{k}{k_{\mathrm{rel}}}
\right)^{\alpha_{\ell}}
\right].
\end{equation}

The characteristic RDCC signatures are:
\begin{itemize}
\item monopole suppression for $k > k_{\mathrm{rel}}$,
\item quadrupole turnover at $k \sim k_{\mathrm{rel}}$,
\item hexadecapole smoothing for all $k > k_{\mathrm{rel}}$.
\end{itemize}

% ---------------------------------------------------------
\subsection*{A.9 Summary}

The analytical approximations in this appendix provide:
\begin{itemize}
\item compact expressions for RDCC-modified RSD observables,
\item analytical control over the Kaiser and FoG regimes,
\item fast semi-analytic tools for RSD modeling,
\item transparent links between relaxation dynamics and RSD signatures.
\end{itemize}

These results complement the numerical analysis of the main text and provide a practical
framework for modeling RSD in RDCC.

% ---------------------------------------------------------
% APPENDIX B — IMPLEMENTATION OF RDCC RSD IN CLASS
% ---------------------------------------------------------
\section*{Appendix B: Implementation of RDCC RSD in CLASS}
\addcontentsline{toc}{section}{Appendix B: Implementation of RDCC RSD in CLASS}

This appendix describes the implementation of the RDCC-modified growth rate, velocity
divergence, density–velocity coupling, and Fingers-of-God (FoG) damping in the CLASS
Boltzmann code. The modifications extend the RDCC growth and velocity framework of
Companions~XVI–XIX and incorporate the analytical approximations derived in Appendix~A
of this Companion.

The goal is to compute RDCC-modified redshift-space distortions (RSD) and make them
available for DESI, Euclid, SKA, and future spectroscopic surveys.

% ---------------------------------------------------------
\subsection*{B.1 Required CLASS Modules}

RDCC RSD requires modifications in the following CLASS modules:

\begin{itemize}
\item \texttt{background.c} — access to $k_{\mathrm{rel}}(a)$ and RDCC growth functions,
\item \texttt{perturbations.c} — RDCC-modified density and velocity transfer functions,
\item \texttt{nonlinear.c} — RDCC-modified matter and velocity power spectra,
\item \texttt{lensing.c} — anisotropic kernels for RSD multipoles,
\item \texttt{input.c} — new RDCC RSD parameters,
\item \texttt{common.h} — RDCC RSD data structures,
\item \texttt{output.c} — RDCC RSD multipole outputs.
\end{itemize}

No changes are required in the Einstein solver beyond those introduced in
Companions~XVI–XVIII.

% ---------------------------------------------------------
\subsection*{B.2 Input Parameters}

The following new input parameters must be added to \texttt{input.c}:

\begin{itemize}
\item \texttt{rdcc\_rsd = yes/no} — activates RDCC RSD module,
\item \texttt{k\_rel} — relaxation scale,
\item \texttt{alpha\_rel} — suppression exponent,
\item \texttt{chi\_f} — RDCC growth-rate amplitude,
\item \texttt{chi\_delta} — RDCC density suppression amplitude,
\item \texttt{chi\_theta} — RDCC velocity-divergence suppression amplitude,
\item \texttt{chi\_deltatheta} — RDCC density–velocity coupling amplitude,
\item \texttt{chi\_sigma} — RDCC velocity-dispersion amplitude.
\end{itemize}

These parameters mirror the analytical forms introduced in Appendix~A.

% ---------------------------------------------------------
\subsection*{B.3 Modifications to the Growth Rate}

The growth rate is computed in \texttt{perturbations.c} as


\[
f = \frac{d\ln D}{d\ln a}.
\]



In RDCC:
\begin{equation}
f_{\mathrm{RDCC}}(k,z)
=
f(k,z)
\left[
1
-
\chi_{f}
\left(
\frac{k}{k_{\mathrm{rel}}}
\right)^{\alpha_{f}}
\right].
\end{equation}

This modification is applied before computing the velocity transfer functions.

% ---------------------------------------------------------
\subsection*{B.4 Modifications to Density and Velocity Transfer Functions}

The density transfer function is modified as
\begin{equation}
T_{\delta,\mathrm{RDCC}}(k,z)
=
T_{\delta}(k,z)
\left[
1
-
\chi_{\delta}
\left(
\frac{k}{k_{\mathrm{rel}}}
\right)^{\alpha_{\delta}}
\right].
\end{equation}

The velocity divergence transfer function becomes
\begin{equation}
T_{\theta,\mathrm{RDCC}}(k,z)
=
T_{\theta}(k,z)
\left[
1
-
\chi_{\theta}
\left(
\frac{k}{k_{\mathrm{rel}}}
\right)^{\alpha_{\theta}}
\right].
\end{equation}

These modifications propagate into all RSD observables.

% ---------------------------------------------------------
\subsection*{B.5 Density–Velocity Cross Spectrum}

The cross-spectrum is computed in \texttt{perturbations.c} as


\[
P_{\delta\theta}(k,z) = -aHf\,P_{\delta\delta}(k,z).
\]



In RDCC:
\begin{equation}
P_{\delta\theta,\mathrm{RDCC}}(k,z)
=
P_{\delta\theta}(k,z)
\left[
1
-
\chi_{\delta\theta}
\left(
\frac{k}{k_{\mathrm{rel}}}
\right)^{\alpha_{\delta\theta}}
\right].
\end{equation}

This modification affects the anisotropic structure of the redshift-space power spectrum.

% ---------------------------------------------------------
\subsection*{B.6 RDCC Fingers-of-God Implementation}

The standard FoG factor is


\[
\mathrm{FoG}(k,\mu) = \exp[-(k\mu\sigma_{v})^{2}].
\]



In RDCC:
\begin{equation}
\sigma_{v,\mathrm{RDCC}}^{2}
=
\sigma_{v}^{2}
\left[
1
-
\chi_{\sigma}
\left(
\frac{k_{\mathrm{NL}}}{k_{\mathrm{rel}}}
\right)^{\alpha_{\sigma}}
\right],
\end{equation}
and
\begin{equation}
\mathrm{FoG}_{\mathrm{RDCC}}(k,\mu)
=
\exp[-(k\mu\sigma_{v,\mathrm{RDCC}})^{2}].
\end{equation}

This predicts weaker small-scale damping.

% ---------------------------------------------------------
\subsection*{B.7 RDCC Redshift-Space Power Spectrum}

The full RDCC redshift-space power spectrum is implemented as
\begin{equation}
P_{s,\mathrm{RDCC}}(k,\mu,z)
=
\left[
b + f_{\mathrm{RDCC}}(k,z)\mu^{2}
\right]^{2}
P_{\delta\delta,\mathrm{RDCC}}(k,z)
\exp[-(k\mu\sigma_{v,\mathrm{RDCC}})^{2}].
\end{equation}

This expression is evaluated for each $\mu$-bin and $k$-mode.

% ---------------------------------------------------------
\subsection*{B.8 RDCC Multipole Outputs}

The RDCC multipoles are written to:

\begin{itemize}
\item \texttt{cl\_rsd\_P0.dat},
\item \texttt{cl\_rsd\_P2.dat},
\item \texttt{cl\_rsd\_P4.dat}.
\end{itemize}

These files can be used directly for DESI, Euclid, and SKA likelihood analyses.

% ---------------------------------------------------------
\subsection*{B.9 Summary}

The CLASS implementation of RDCC RSD requires:

\begin{itemize}
\item RDCC-modified growth rate and density transfer functions,
\item RDCC-modified velocity divergence and cross spectra,
\item RDCC-modified Fingers-of-God damping,
\item new CLASS parameters and data structures,
\item RDCC multipole outputs for DESI, Euclid, and SKA.
\end{itemize}

This implementation provides a complete numerical framework for computing RDCC RSD
observables and enables direct comparison with current and future spectroscopic surveys.

% ---------------------------------------------------------
% APPENDIX C — NUMERICAL SETTINGS AND VALIDATION TESTS
% ---------------------------------------------------------
\section*{Appendix C: Numerical Settings and Validation Tests}
\addcontentsline{toc}{section}{Appendix C: Numerical Settings and Validation Tests}

This appendix summarizes the numerical settings, convergence tests, and validation procedures
used to compute the RDCC redshift-space distortion (RSD) observables presented in this
Companion. The goal is to ensure reproducibility and to provide guidance for implementing
RDCC RSD in CLASS and related analysis pipelines.

% ---------------------------------------------------------
\subsection*{C.1 Numerical Resolution Requirements}

Accurate computation of RDCC-modified RSD requires sufficient resolution in $k$-space,
$\mu$-sampling, and redshift.

Recommended settings:
\begin{itemize}
\item $k_{\mathrm{max}} \geq 0.8\,h\,\mathrm{Mpc}^{-1}$ for RSD multipoles,
\item $\Delta k / k \leq 0.03$ for stable RDCC suppression factors,
\item $N_{\mu} \geq 200$ for anisotropic power-spectrum integration,
\item $\Delta z \leq 0.01$ for growth-rate evolution,
\item $\ell_{\mathrm{max}} \geq 6$ for RSD multipoles ($\ell = 0,2,4$).
\end{itemize}

These settings ensure convergence of the RDCC-modified RSD spectra at the percent level.

% ---------------------------------------------------------
\subsection*{C.2 Convergence Tests}

We perform three classes of convergence tests:

\paragraph{(1) $k$-space resolution.}
The RDCC suppression factor


\[
1 - \chi\left(\frac{k}{k_{\mathrm{rel}}}\right)^{\alpha}
\]


requires fine sampling near $k \sim k_{\mathrm{rel}}$.
Convergence is achieved when:


\[
\frac{\Delta P_{\ell}^{\mathrm{RDCC}}}{P_{\ell}^{\mathrm{RDCC}}} < 0.5\%
\quad \text{for } \ell = 0,2,4.
\]



\paragraph{(2) $\mu$-sampling.}
The anisotropic structure of $P_{s}(k,\mu)$ is sensitive to RDCC velocity divergence.
Convergence is achieved when:


\[
N_{\mu} \geq 200
\quad \Rightarrow \quad
\frac{\Delta P_{\ell}}{P_{\ell}} < 0.3\%.
\]



\paragraph{(3) Growth-rate evolution.}
The scale-dependent growth rate $f_{\mathrm{RDCC}}(k,z)$ requires fine redshift sampling.
Convergence is achieved when:


\[
\Delta z \leq 0.01
\quad \Rightarrow \quad
\frac{\Delta f_{\mathrm{RDCC}}}{f_{\mathrm{RDCC}}} < 0.2\%.
\]



% ---------------------------------------------------------
\subsection*{C.3 Validation Against $\Lambda$CDM}

RDCC RSD must reduce to $\Lambda$CDM in the limit:


\[
\chi_{f},\chi_{\delta},\chi_{\theta},\chi_{\delta\theta},\chi_{\sigma} \rightarrow 0,
\qquad
k_{\mathrm{rel}} \rightarrow \infty.
\]



We validate:
\begin{itemize}
\item $P_{0,\mathrm{RDCC}} \rightarrow P_{0}$ at the $<0.1\%$ level,
\item $P_{2,\mathrm{RDCC}} \rightarrow P_{2}$ at the $<0.1\%$ level,
\item $P_{4,\mathrm{RDCC}} \rightarrow P_{4}$ at the $<0.2\%$ level,
\item $f_{\mathrm{RDCC}} \rightarrow f$ at the $<0.1\%$ level,
\item $\sigma_{v,\mathrm{RDCC}} \rightarrow \sigma_{v}$ at the $<0.1\%$ level.
\end{itemize}

These checks ensure that RDCC is a controlled extension of $\Lambda$CDM.

% ---------------------------------------------------------
\subsection*{C.4 Validation of RDCC RSD Signatures}

We validate the characteristic RDCC signatures:

\paragraph{(1) Monopole suppression.}


\[
P_{0,\mathrm{RDCC}}(k)
<
P_{0}(k)
\quad \text{for } k > k_{\mathrm{rel}}.
\]



\paragraph{(2) Quadrupole turnover.}


\[
P_{2,\mathrm{RDCC}}(k)
>
P_{2}(k)
\quad \text{for } k < k_{\mathrm{rel}},
\]


and


\[
P_{2,\mathrm{RDCC}}(k)
<
P_{2}(k)
\quad \text{for } k > k_{\mathrm{rel}}.
\]



\paragraph{(3) Hexadecapole smoothing.}


\[
P_{4,\mathrm{RDCC}}(k)
<
P_{4}(k)
\quad \text{for all } k > k_{\mathrm{rel}}.
\]



\paragraph{(4) Reduced velocity dispersion.}


\[
\sigma_{v,\mathrm{RDCC}} < \sigma_{v}.
\]



\paragraph{(5) Weaker FoG damping.}


\[
\mathrm{FoG}_{\mathrm{RDCC}}(k,\mu)
>
\mathrm{FoG}(k,\mu)
\quad \text{for } k > k_{\mathrm{rel}}.
\]



These signatures provide a clean observational handle on the relaxation scale.

% ---------------------------------------------------------
\subsection*{C.5 Numerical Stability and Error Control}

To ensure numerical stability:

\begin{itemize}
\item enforce positivity of all RDCC-modified spectra,
\item apply spline smoothing to RDCC suppression factors,
\item ensure smooth transitions in $f_{\mathrm{RDCC}}(k,z)$,
\item use adaptive $\mu$-sampling near $\mu \approx 1$,
\item apply exponential damping for $k \gg k_{\mathrm{rel}}$.
\end{itemize}

These steps prevent numerical artifacts in the RDCC RSD multipoles.

% ---------------------------------------------------------
\subsection*{C.6 Summary}

This appendix provides:
\begin{itemize}
\item recommended numerical settings for RDCC RSD,
\item convergence tests for $k$-space, $\mu$-sampling, and growth-rate evolution,
\item validation against $\Lambda$CDM,
\item verification of RDCC RSD signatures,
\item guidelines for stable numerical implementation.
\end{itemize}

These procedures ensure that RDCC RSD predictions are accurate, reproducible, and
consistent with the analytical and numerical framework developed in this Companion.


\end{document}
